% LSPV-2 REPORT::SUCTYPE1::2026-10-12T18:21:16Z
\documentclass[11pt]{article}
\usepackage[margin=1in]{geometry}
\usepackage{amsmath,amssymb,booktabs,array,hyperref}
\usepackage[T1]{fontenc}
\hypersetup{colorlinks=true,linkcolor=blue,urlcolor=blue}
\newcommand{\SAT}{\mathsf{3SAT}_{\rm exp}}
\newcommand{\SSAT}{\mathsf{Succinct\mbox{-}3SAT}}
\newcommand{\OPEN}{\textbf{OPEN}}
\newcommand{\EXACT}{\textbf{PROVED/EXACT}}
\newcommand{\KNOWN}{\textbf{KNOWN THEOREM}}
\newcommand{\COMP}{\textbf{COMPUTATIONAL}}
\newcommand{\RET}{\textbf{FALSIFIED/RETIRED}}

\title{LSPV-2 SUCTYPE1: the succinct-tableau type firewall}
\author{UNIVERSAL\_CORE\_TESTS (LSPV-2)}
\date{2026-10-13 03:21 KST}

\begin{document}
\maketitle

\section*{Executive status}
\textbf{Result.}  A succinct circuit descriptor is not an input to an ordinary
explicit-input SAT decider.  There are exactly two direct ways to connect the
objects: materialize the represented formula, paying its output size, or build
a new lazy-oracle simulation and charge every descriptor evaluation.  Treating
the descriptor itself as a SAT instance instead defines succinct SAT, an
NEXP-complete language.  This closes a representation loophole in the current
encoding-scale firewall but gives no lower bound for ordinary explicit 3-SAT.

\begin{center}
\begin{tabular}{>{\raggedright\arraybackslash}p{0.42\linewidth}p{0.18\linewidth}p{0.30\linewidth}}
\toprule
Claim & Status & Consequence\\
\midrule
SUCCINCT-TABLEAU-TYPE-FIREWALL1 & \EXACT{} & Type and all bridge costs are explicit.\\
Succinct 3-SAT is NEXP-complete & \KNOWN{} & A descriptor solver is not an ordinary SAT solver.\\
SUCCINCT-DESCRIPTOR-AS-ORDINARY-SAT-CONTRACTION1 & \RET{} & It silently changes the input language/model.\\
Universal clean-core transfer to every polynomial-time SAT decider & \OPEN{} & No extraction or model-matched contradiction is known.\\
P versus NP & \OPEN{} & This report proves neither separation nor equality.\\
\bottomrule
\end{tabular}
\end{center}

\section*{Predecessor write-recovery preflight}
The immediately preceding LSPV-1 run was verified directly from GitHub at
\texttt{research/agents/lspv-1/2026-10-13\_0311\_LSPV-1\_HIERFIRE1.tex}.
It is substantive and contains both the required LSPV-1 header and final
sentinel.  No recovery prompt or repair was required.

\section*{Barrier-transfer literature preflight}
The two active obstacles were: (i) an input of length $s$ describing an object
of length $L$ had been compared to an algorithm whose input convention is the
$L$-bit explicit object; and (ii) lazy access had not been separated from a
direct call to that explicit-input algorithm.  Primary literature was checked
before the analysis.  Papadimitriou and Yannakakis prove the representation
shift for succinct encodings and, in particular, the NEXP phenomenon for
succinct versions of NP-complete problems \cite{PY1986}.  Cook's
nondeterministic time hierarchy supplies the unconditional class-separation
context \cite{Cook1973}.  Cook--Levin supplies the explicit tableau baseline
\cite{Cook1971}.

\begin{itemize}
\item \textbf{Directly transferable theorem:} a circuit-succinct version of
3-SAT is NEXP-complete under the standard succinct representation.  Object
correspondence: descriptor circuit $D$ $\leftrightarrow$ the bits of an
explicit formula $y_D$.  This theorem concerns the descriptor language, not
ordinary explicit 3-SAT.
\item \textbf{Useful structural analogy requiring a new transfer theorem:}
compressed-data random access and VLSI decompression emphasize access and
routing costs, but neither proves a lower bound for an unrestricted uniform
SAT machine.
\item \textbf{Misleading analogy / failed transfer:} ``the descriptor is
short, so run the polynomial-time SAT decider on it.''  The descriptor is not a
well-typed input to that decider.
\end{itemize}

\section*{Claim card and exact quantifiers}
Fix a binary, self-delimiting, explicit encoding of 3-CNF formulas and let
$\SAT\subseteq\{0,1\}^*$ denote the resulting language.  A descriptor is a
Boolean circuit $D$ with $a$ address inputs and one output.  Together with a
length field $L\leq 2^a$, it denotes
\[
       y_D=D(0)D(1)\cdots D(L-1)\in\{0,1\}^L.
\]
Malformed strings are rejected by the explicit syntax checker.  Let
$A$ be any deterministic uniform machine that decides $\SAT$ and whose
worst-case time on length-$N$ explicit inputs is $t_A(N)$; under the conditional
assumption SAT$\in$P, $t_A(N)\leq N^{K_A+o(1)}$ for some algorithm-dependent
constant $K_A\geq 1$.

\paragraph{Target tested.}
For every such $A$, can a circuit descriptor of a tableau formula be fed to
$A$ at cost polynomial in $|D|$, thereby avoiding the explicit encoding scale?
The answer is no for a direct call; any repair is one of the two charged
bridges below or a different language.

\section*{SUCCINCT-TABLEAU-TYPE-FIREWALL1}
\textbf{Theorem (\EXACT{}, pending independent review).}
For every descriptor $(D,L)$ and every explicit-input decider $A$ above:
\begin{enumerate}
\item $A(D)$ has no reason to equal $[y_D\in\SAT]$.  It applies $A$ to the
literal encoding of the circuit, not to the represented $L$-bit formula.
\item An eager bridge that writes $y_D$ before invoking $A$ costs at least
$\Omega(L)$ bit writes in a bit-cost multitape model, or
$\Omega(L/\log L)$ word writes on a word RAM with word size
$\Theta(\log L)$, in addition to descriptor evaluation and $t_A(L)$.
\item A lazy bridge must define a new composed machine.  If its simulation
makes $q$ logical input-bit probes and one evaluation of $D(i)$ costs
$\tau_D(i)$, then the access charge is exactly
$\sum_{j=1}^{q}\tau_D(i_j)$, plus the simulated machine's remaining work.
There is no universal theorem here forcing $q=L$, nor one bounding the full
simulation by a polynomial in $|D|$.
\item A machine that accepts $(D,L)$ exactly when $y_D\in\SAT$ decides
$\SSAT$, not $\SAT$.  For the standard circuit representation this language
is NEXP-complete \cite{PY1986}.  By nondeterministic time hierarchy it is not
in P.  The assumption P$=$NP does not turn this descriptor problem into an
ordinary polynomial-time problem.
\end{enumerate}

\paragraph{Proof.}
Item 1 is the definition of a language as a set of strings: $D$ and $y_D$ are
different strings with different parsers.  In item 2, the output interface must
emit $L$ bits (or $\lceil L/w\rceil$ words), so the stated write lower bound is
model-matched and independent of the formula semantics.  Item 3 is exact
accounting for the oracle substitution $y_D[i]\mapsto D(i)$; it deliberately
makes no all-algorithm probe lower bound.  Item 4 is the standard succinct
representation theorem together with the nondeterministic hierarchy theorem.
None of the four steps uses P$\ne$NP. \hfill$\square$

For a $T$-step Cook--Levin-style tableau whose explicit formula length is
$L=T^{1+o(1)}$, eager use of a hypothetical polynomial-time SAT decider costs
\[
  T^{1+o(1)}+t_A(T^{1+o(1)})=T^{K_A+o(1)}
\]
when $K_A\geq1$.  Thus it does not evade the current HIERFIRE1 inequality
$\max\{\rho,\lambda K_A\}<1$ with standard $\lambda=1$.  This is an input/output
accounting conclusion, not a SAT time lower bound.

\section*{Bounded reproducible hostile-control audit}
The deterministic script
\texttt{experiments/review/succinct\_tableau\_type\_firewall\_audit.py}
constructed explicit fixed-width encodings and random-access descriptors for
eight families at $n\in\{4,8,16,32,64,128,256\}$:
2SAT, Horn, affine/Tseitin, bounded treewidth, independent products, padding,
redundant clauses, and an alternate encoding.  It audited six equivalent
wrappers (DIRECT, COPY, DEAD, DELAY, INTERP, RECOMP), eight arbitrary polynomial
exponents $K=1,\dots,8$, and separated descriptor construction, materialization
bit/word writes, lazy decoding, extraction, reconstruction, and one-bit output.
For $n\leq8$ it also exhaustively checked the represented formula semantics.

The run completed in $3.2$ seconds as one sub-30-second chunk: 1,484 exact
checks, 56 family-size rows, 448 exponent ledgers, 336 wrapper rows, and zero
failures.  The largest observed finite compression ratio was $723.2$: an
85-bit descriptor generated a 61,472-bit padding formula at $n=256$.  That
formula is a conjunction of tautologies and is easy.  Hence a large
description/expansion ratio is not a hardness witness.  At $n=8$ the exact
model counts were respectively $9,9,4,9,81,256,128,81$ for the eight families,
while all decision outputs remained one bit.

Result file:
\texttt{results/raw/review/2026-10-13\_succinct\_tableau\_type\_firewall.json}.
Script SHA-256:
\texttt{4ebebcbac46c605421c5c9e3cabeb0609c2cb0023ca3a29aaa61fb43142e29d4}.
Result SHA-256:
\texttt{d71a500e686991afadba1cdb26d745211fb9351a38b1b49e087a68278d666d8c}.
The audit is \COMP{} and does not establish an asymptotic lower bound.

\section*{SAME\_REALIZATION and complete cost ledger}
The descriptor circuit, its expanded formula, and a run of $A$ are three
distinct realizations.  They may be equated only after one of these interfaces
is specified and charged:
\[
\begin{array}{ll}
\text{eager:}& C_{\rm build}(D)+C_{\rm eval}(D,L)+C_{\rm write}(L)+t_A(L),\\
\text{lazy:}& C_{\rm build}(D)+\sum_{j\le q}\tau_D(i_j)+C_{\rm sim}(A,q),\\
\text{direct succinct:}& t_B(|D|+\log L)\text{ for a different algorithm }B.
\end{array}
\]
Extraction and reconstruction are zero in the decision-only controls; the final
output is one bit.  No semantic-state count is converted into runtime.  COPY,
DEAD, DELAY, INTERP, and RECOMP wrappers preserve the answer while changing
internal traces and costs, so no unqualified exposure theorem can be invariant
over implementations.

\section*{Relativization, algebrization, Natural Proofs, and nonuniformity}
The type distinction and output-write charge relativize and survive arbitrary
oracles; they therefore do not bypass the Baker--Gill--Solovay barrier.  They
are not an algebraic lower-bound argument and do not evade algebrization.  No
large constructive circuit property is proposed, so the Natural Proofs
framework is not a positive route here.  A nonuniform family of descriptors
also does not yield a uniform ordinary-input SAT algorithm: a uniform
descriptor constructor and every access/compilation cost must still be charged.
The result is a firewall against a false shortcut, not a nonrelativizing
separation technique.

\section*{Post-quantum technique ledger}
\begin{itemize}
\item Schur squares, quadratic-form relation spaces, Hilbert functions, and
syzygies: \textbf{FAILED TRANSFER} for this input-type obstacle; there is no
code multiplication map in the theorem.
\item Lossiness/collision entropy: \textbf{ANALOGY/CANDIDATE}.  A short
descriptor can have a huge expansion fiber, but the padding control shows that
collision/compression alone can expose easy structure rather than hardness.
No bridge to an arbitrary SAT decider is known.
\item Worst-case/average-case and external extraction (LWE lesson):
\textbf{ANALOGY/CANDIDATE}.  Black-box extraction could avoid claims about an
algorithm's internals, but no worst-case/average-case reduction or extractor
for this clean core is known.
\item Special soundness/collision extraction and hash/random-oracle methods:
\textbf{FAILED TRANSFER} absent a construction-specific unconditional witness.
Cryptographic assumptions are inadmissible for the target separation.
\end{itemize}

\section*{CORPUS REVIEWED}
The GitHub tree was surveyed before substantive work.  Reviewed sources include
root \texttt{main.tex}, \texttt{AGENTS.md}, \texttt{CODEX.md}, and the complete
\texttt{.codex/skills/pnp-research/SKILL.md} protocol and references; all task
specifications under \texttt{tasks/}; actual \texttt{main4.tex} and
\texttt{main5.tex}; the newest verified MAIN
\texttt{2026-10-13\_0148\_MAIN\_TRACECONJ1.tex}; the materially relevant
historical MAIN \texttt{SYNTHESIS-0024}; newest verified LSPV-1 HIERFIRE1,
LSPV-2 INFOFLOWCHAIN1, LSPV-3 REVUNCOMP1, and LSPV-4 READCEIL1; historical
LSPV-2 RECURCLOSE1; the cost-certificate spine; the HEC, VSC-HEC, G-PPC-HEC,
CRS-PPC, CCSPC, and search-spine explorations; relevant governance notes; and
the experiment/result artifacts cited by active claims.  Repository searches
covered ``succinct tableau'', ``succinct SAT'', ``reduction image'',
``descriptor'', and the retired resource routes.

\section*{DUPLICATION / CONTRADICTION AUDIT}
The corpus already said qualitatively that succinct tableaux change the
problem and require compilation costs.  The new contribution is the exact
ordinary-input type statement, the exhaustive eager/lazy/direct trichotomy,
the bit- and word-write charges, the NEXP class ledger, and a hostile-control
experiment.  It does not duplicate LOG-SUCCINCT-RESOURCE, which concerned
resource equivalence, or RECURCLOSE1, which concerned recursive closure over a
reduction image.  No verified theorem was contradicted.  The new result closes
only the succinct-representation loophole in HIERFIRE1.

\section*{WHAT CHANGED RELATIVE TO ROOT MAIN}
The root LSPV-2 region previously ended at INFOFLOWCHAIN1 and listed succinct
verification only as a caution.  It is reconciled to add SUCTYPE1 as an exact
type/cost firewall, to retire the descriptor-as-ordinary-SAT shortcut, and to
retain the universal transfer theorem as OPEN.  No executive or other
specialist-owned text is changed by LSPV-2.

\section*{ADOPTED / REJECTED / RETIRED PRIOR COMPONENTS}
\begin{itemize}
\item \textbf{ADOPTED:} HIERFIRE1's exact
$T^{\max\{\rho,\lambda K_A\}+o(1)}$ accounting; INFOFLOWCHAIN1's rejection of
state-count-to-time inference; the one-bit decision-output control; and the
wrapper-invariance requirement.
\item \textbf{REJECTED:} any premise that every algorithm must materialize a
chosen compatibility object or read all represented bits.
\item \textbf{RETIRED:} SUCCINCT-DESCRIPTOR-AS-ORDINARY-SAT-CONTRACTION1 and
any uncharged substitution of $|D|$ for the explicit input length $L$.
\end{itemize}

\section*{STALE OR OVERSTATED CLAIMS FOUND}
Any prose reading ``a succinct tableau gives a short SAT instance'' without
naming the descriptor language is overstated.  The short string is an instance
of a succinct language; expanding it reinstates the explicit length, and lazy
access defines a new composed computation.  The root MAIN already warned about
this distinction, so the repair is a formal strengthening rather than a
reversal.  No current verified claim that P$\ne$NP was found.

\section*{Faithful comparison to main4/main5 and handoff}
The actual \texttt{main4.tex} and \texttt{main5.tex} describe compatibility
maps, semantic histories, reconstruction, and universal-cost ambitions; they
do not contain a theorem making a descriptor a well-typed input to every SAT
decider.  This report does not redefine their clean-core objects.  It removes
one proposed bridge and leaves the same load-bearing target:
\[
\boxed{\text{Find a representation-invariant, all-algorithm transfer theorem
with a model-matched unconditional lower bound.}}
\]
The best next test is a special-image promise formulation: quantify exactly
over algorithms correct on all explicit formulas, restrict their inputs to a
uniform tableau image, and ask for an externally observable resource whose
extraction cost includes image recognition and all descriptor/expansion work.
Until that arrow is proved, P versus NP remains open.

\begin{thebibliography}{9}
\bibitem{Cook1971}
S. A. Cook, ``The complexity of theorem-proving procedures,''
\emph{Proceedings of STOC}, 1971, pp. 151--158.
\bibitem{Cook1973}
S. A. Cook, ``A hierarchy for nondeterministic time complexity,''
\emph{Journal of Computer and System Sciences} 7(4), 1973, pp. 343--353,
doi:10.1016/S0022-0000(73)80028-5.
\bibitem{PY1986}
C. H. Papadimitriou and M. Yannakakis, ``A note on succinct representations
of graphs,'' \emph{Information and Control} 71(3), 1986, pp. 181--185,
doi:10.1016/S0019-9958(86)80009-2.
\end{thebibliography}

\end{document}
% END_OF_VERIFIED_REPORT::LSPV-2::SUCTYPE1
